\documentclass[11pt,a4paper]{article}
\usepackage[utf8]{inputenc}
\usepackage[T1]{fontenc}
\usepackage[bahasa]{babel}
\usepackage[margin=2.5cm]{geometry}
\usepackage{xcolor}
\usepackage{mdframed}
\usepackage{enumitem}
\usepackage{hyperref}
\hypersetup{colorlinks=true, urlcolor=blue, linkcolor=black}

% Gaya kotak komentar reviewer & jawaban penulis
\definecolor{revbg}{HTML}{EEF3F8}
\definecolor{ansac}{HTML}{2E7D32}
\newmdenv[backgroundcolor=revbg, linecolor=revbg, skipabove=6pt, skipbelow=4pt,
  innertopmargin=6pt, innerbottommargin=6pt]{revbox}
\newcommand{\reviewer}[1]{\vspace{4pt}\noindent\textbf{\textcolor{black}{#1}}}
\newcommand{\authors}[1]{\vspace{2pt}\noindent\textcolor{ansac}{\textbf{Tanggapan penulis:}} #1\par}
\newcommand{\loc}[1]{\vspace{1pt}\noindent{\small\textit{Lokasi perubahan: #1}}\par\vspace{6pt}}

\title{\textbf{Tanggapan terhadap Reviewer}\\[4pt]
\large Naskah ID \#42551\\
\normalsize ``Design of a cloud backup system based on multi-compression algorithms
for ransomware attack mitigation''}
\author{Hero Yudo Martono, Nafisah Rahmadani Harahap\\
\small Politeknik Elektronika Negeri Surabaya, Indonesia}
\date{\today}

\begin{document}
\maketitle

\noindent Kami mengucapkan terima kasih kepada Editor dan seluruh reviewer atas masukan yang
membangun dan terperinci. Kami telah merevisi naskah secara saksama untuk menanggapi setiap
komentar. Berikut tanggapan kami secara poin-per-poin. Komentar reviewer ditampilkan dalam kotak
berwarna; tanggapan kami beserta lokasi perubahan mengikuti setiap komentar.

\medskip
\noindent\textbf{Ringkasan revisi utama:}
\begin{itemize}[leftmargin=1.4em, itemsep=1pt]
  \item Memperjelas ketidaksesuaian ukuran dataset: dataset terdiri atas \textbf{43 berkas} secara
        konsisten di seluruh naskah (nilai ``113'' sebelumnya telah dikoreksi).
  \item Abstrak kini menyatakan secara eksplisit bahwa serangan ransomware bersifat
        \textbf{disimulasikan dan terkontrol} (tanpa mengeksekusi malware nyata) pada
        \textbf{dataset sintetis berukuran terbatas}.
  \item Menambahkan pembahasan yang menjelaskan bahwa FPR/FNR 0\% adalah konsekuensi logis
        dari deteksi berbasis integritas \textit{hash} pada lingkungan terkontrol, bukan klaim umum.
  \item Mengoreksi urutan teknis menjadi \textbf{compress-then-encrypt} dan menjalankan ulang seluruh eksperimen.
  \item Menambahkan paragraf State-of-the-Art, paragraf \textit{research gap} dengan formulasi
        entropi/kompresi, dan perbandingan dengan literatur berorientasi deteksi.
  \item Memvalidasi sistem pada \textbf{dua platform} (Windows dan Linux) di AWS EC2, disertai tabel
        konsistensi lintas-platform.
  \item Menambahkan \textbf{analisis statistik} (5 uji-berulang): rata-rata~$\pm$~standar deviasi
        waktu kompresi, \textit{throughput}, dan utilisasi CPU/RAM.
  \item Memperluas arah penelitian mendatang melampaui Random Forest.
\end{itemize}

% =====================================================================
\section*{Tanggapan untuk Editor}
% =====================================================================

\reviewer{Komentar Editor 1. Perkuat struktur, kejelasan, dan argumentasi; mulai dengan pernyataan
tesis yang kuat; tiap paragraf memiliki kalimat topik dan alur logis.}
\begin{revbox}\small
Perbaiki struktur dan argumentasi secara keseluruhan; mulai pendahuluan dengan tesis yang jelas
dan didukung bukti yang kuat.
\end{revbox}
\authors{Bagian Pendahuluan telah disusun ulang. Kini dibuka dengan ancaman ransomware yang
meningkat (didukung statistik kuantitatif), membangun rantai logis dari masalah (cadangan
terhubung yang rentan, kompresi algoritma tunggal) menuju \textit{research gap}, dan ditutup
dengan pernyataan kontribusi yang eksplisit.}
\loc{Bagian~1 (Pendahuluan), paragraf yang direvisi dan paragraf State-of-the-Art baru.}

\reviewer{Komentar Editor 2. Tonjolkan \textit{state of the art} di Pendahuluan; bandingkan temuan
dengan hasil terdahulu pada Hasil \& Pembahasan; definisikan kontribusi secara eksplisit.}
\begin{revbox}\small
Tekankan state of the art; bandingkan dengan hasil sebelumnya; nyatakan kebaruan dengan jelas.
\end{revbox}
\authors{Kami menambahkan paragraf State-of-the-Art khusus yang merangkum riset deteksi dan
mitigasi ransomware terkini (\textit{ensemble} ML $>$99\%, \textit{deep learning} 96--99\%,
detektor waktu-nyata), dan secara eksplisit memposisikan kontribusi kami pada lapisan
\emph{mitigasi/pemulihan}, bukan bersaing pada akurasi deteksi. Hasil kini dibandingkan dengan
literatur berorientasi deteksi.}
\loc{Bagian~1 (paragraf State-of-the-Art) dan Bagian~2 (Penelitian Terkait, paragraf perbandingan baru).}

\reviewer{Komentar Editor 3. Bagian Metode harus mencakup pendekatan standar dan baru disertai
justifikasi validitas, serta prosedur langkah-demi-langkah yang dapat direproduksi.}
\begin{revbox}\small
Metode harus dapat direproduksi dan menjustifikasi validitas pendekatan.
\end{revbox}
\authors{Bagian Metode kini menjelaskan alur \emph{compress-then-encrypt} yang telah dikoreksi
beserta justifikasinya (data terenkripsi ber-entropi tinggi tidak dapat dikompresi). Seluruh
langkah (garis dasar SHA-256, kompresi, enkripsi AES-256, penyimpanan \textit{air-gapped},
deteksi, pemulihan) dijelaskan berurutan, dan implementasi lengkap dijalankan pada AWS EC2
sehingga prosedurnya dapat direproduksi.}
\loc{Bagian~3 (Metode), subbagian arsitektur, enkripsi, kompresi, serta cadangan/pemulihan.}

\reviewer{Komentar Editor 4. Naskah kurang pembahasan kritis, perbandingan, dan interpretasi implikasi.}
\begin{revbox}\small
Tambahkan pembahasan kritis: apa implikasinya, dan apa yang berguna ke depan?
\end{revbox}
\authors{Kami menambahkan pembahasan kritis atas hasil FPR/FNR (makna dan keterbatasannya),
analisis konsistensi lintas-platform, dan Kesimpulan yang diperluas secara substansial mencakup
keterbatasan dan arah penelitian mendatang.}
\loc{Bagian~4 (Hasil \& Pembahasan: diskusi FPR/FNR, subbagian lintas-platform) dan Bagian~5 (Kesimpulan).}

% =====================================================================
\section*{Tanggapan untuk Reviewer A}
% =====================================================================

\reviewer{A1. Pertanyaan strategis: kesesuaian ruang lingkup, objektif/pendekatan, kebaruan,
kontribusi terhadap pengetahuan, implikasi masa depan, dan mengapa jurnal harus mempublikasikannya.}
\begin{revbox}\small
Bagaimana naskah sesuai dengan ruang lingkup jurnal dan mengapa menarik? Apa objektif dan
pendekatannya? Apa kontribusi barunya? Bagaimana menambah pengetahuan yang ada? Apa implikasi
masa depannya? Mengapa jurnal ini harus mempublikasikannya?
\end{revbox}
\authors{Naskah sesuai ruang lingkup jurnal (sistem penyimpanan awan yang aman dan efisien).
Objektifnya adalah cadangan awan tahan-ransomware yang terintegrasi dan hemat biaya. Kebaruannya
adalah integrasi isolasi \textit{air-gapped} logis, kompresi multi-algoritma adaptif, dan deteksi
integritas multi-indikator dalam satu kerangka---integrasi yang pada studi terdahulu hanya
ditangani secara terpisah. Poin-poin ini kini dinyatakan eksplisit di Pendahuluan dan Penelitian Terkait.}
\loc{Bagian~1 (paragraf kontribusi) dan Bagian~2 (paragraf kebaruan).}

\reviewer{A2. Abstrak tidak menyatakan bahwa serangan disimulasikan/terkontrol (bukan malware
nyata), serta dataset bersifat sintetis dan terbatas.}
\begin{revbox}\small
Abstrak harus menyatakan serangan disimulasikan dan terkontrol, serta dataset sintetis/terbatas.
\end{revbox}
\authors{Abstrak kini menyatakan secara eksplisit bahwa validasi menggunakan \emph{model serangan
tersimulasi yang terkontrol (tanpa mengeksekusi malware nyata)} pada \emph{dataset sintetis
terbatas berisi 43 berkas}.}
\loc{Abstrak (direvisi).}

\reviewer{A3. Pendahuluan perlu menegaskan bahwa studi terdahulu membahas keamanan cadangan,
mitigasi ransomware, dan kompresi secara terpisah, serta belum menganalisis arsitektur cadangan
tahan-ransomware, penyimpanan \textit{immutable}, \textit{snapshot recovery}, deduplikasi, atau
\textit{zero-trust backup}.}
\begin{revbox}\small
Perjelas research gap terhadap immutable storage, snapshot recovery, deduplication, zero-trust backup.
\end{revbox}
\authors{Paragraf \textit{research gap} ditambahkan, menegaskan bahwa karya terdahulu membahas
aspek-aspek tersebut secara terpisah, dan memposisikan kerangka terintegrasi kami terhadapnya.
Konsep terkait (\textit{immutable}/WORM, \textit{zero-trust backup}, deduplikasi) kini dirujuk dan
dimasukkan ke dalam arah penelitian mendatang.}
\loc{Bagian~1 (paragraf research gap) dan Bagian~5 (arah penelitian mendatang).}

\reviewer{A4. Masalah Metode: (i) alur kerja membingungkan; (ii) kompresi adaptif hanya berbasis
aturan; (iii) pemetaan CVE kurang rigor teknis; (iv) deteksi terlalu bergantung pada perubahan
\textit{hash} dan metadata.}
\begin{revbox}\small
Perjelas alur kerja; nyatakan kompresi berbasis aturan secara eksplisit; perkuat pemetaan CVE;
bahas ketergantungan pada hash/metadata.
\end{revbox}
\authors{(i) Alur kerja ditulis ulang dalam langkah berurutan yang jelas dengan urutan
compress-then-encrypt yang telah dikoreksi. (ii) Pemilihan kompresi kini dijelaskan secara
eksplisit sebagai \emph{berbasis aturan} (menurut tipe dan ukuran berkas). (iii) Setiap serangan
tersimulasi dipetakan ke pola CVE (CVE-2017-0144 \textit{encrypt-and-hold}, CVE-2018-8453
\textit{metadata/header corruption}, CVE-2021-36934 \textit{overwrite-and-corrupt}) disertai
deskripsi modifikasi berkas terkontrol yang dilakukan. (iv) Kami kini secara eksplisit mengakui
bahwa deteksi berbasis \textit{hash}/metadata adalah mekanisme integritas (mendeteksi \emph{bahwa}
berkas berubah) dan membahas keterbatasannya terhadap penghindaran canggih (FPE, \textit{fileless}, LotL).}
\loc{Bagian~3 (Metode: arsitektur, kompresi, simulasi serangan, deteksi) dan Bagian~4 (diskusi FPR/FNR).}

\reviewer{A5. Hasil kurang analisis statistik terperinci (uji berulang, standar deviasi, utilisasi
CPU/RAM, \textit{throughput}, energi, biaya \textit{cloud}) dan perbandingan dengan
deduplikasi/incremental backup; beberapa gambar tidak dirujuk.}
\begin{revbox}\small
Tambahkan analisis statistik dan perbandingan; pastikan semua gambar dirujuk.
\end{revbox}
\authors{Kami memperkuat analisis statistik secara substansial. Dalam subbagian baru, kami kini
melaporkan \textbf{statistik uji-berulang (5 \textit{trials})} untuk setiap algoritma pada seluruh
dataset: rata-rata dan standar deviasi waktu kompresi, \textit{throughput} (MB/s), serta utilisasi
CPU/RAM (diukur dengan \texttt{psutil}) di AWS EC2. Standar deviasi yang sangat kecil (mis. Snappy
$0{,}003$~s, Zstandard $0{,}016$~s) mengonfirmasi kinerja yang stabil, dan rentang \textit{throughput}
(Snappy/LZ4 $>590$~MB/s vs.\ Gzip $30{,}5$~MB/s) memberikan dasar objektif bagi seleksi adaptif. Semua
gambar kini dirujuk di teks sebelum ditampilkan. Selain itu, kami menambahkan \textbf{perbandingan
kuantitatif strategi adaptif yang diusulkan terhadap baseline} (tanpa kompresi, Gzip statis, dan
Zstandard statis) pada seluruh dataset. Perbandingan ini mengungkap \textit{trade-off} kecepatan
vs.\ penyimpanan yang jelas: strategi adaptif sekitar $21\times$ lebih cepat daripada Gzip statis
($0{,}87$~s vs.\ $18{,}84$~s; $611{,}8$ vs.\ $28{,}2$~MB/s) sambil tetap mencapai penghematan
$43{,}3\%$, menunjukkan bahwa keunggulan utamanya adalah keseimbangan optimal antara efisiensi
penyimpanan dan kecepatan proses untuk cadangan ber-RTO rendah. Profil energi dan perbandingan
langsung dengan deduplikasi/incremental backup tetap dicantumkan sebagai arah penelitian mendatang.}
\loc{Bagian~4 (subbagian analisis statistik dan subbagian perbandingan baseline baru dengan tabel; seluruh rujukan gambar) dan Bagian~5 (sisa arah penelitian mendatang).}

\reviewer{A6. Arah penelitian mendatang pada Kesimpulan terlalu sempit (hanya Random Forest).}
\begin{revbox}\small
Perluas arah penelitian mendatang melampaui Random Forest.
\end{revbox}
\authors{Arah penelitian mendatang kini mencakup: seleksi kompresi berbasis \textit{machine
learning}, pengujian dengan ransomware nyata pada \textit{sandbox} yang aman, integrasi cadangan
\textit{immutable}/WORM dan \textit{zero-trust}, analisis keamanan formal dan manajemen kunci,
evaluasi skala perusahaan (\textit{throughput}, CPU/RAM, energi, biaya \textit{cloud}), serta
perbandingan langsung dengan \textit{incremental}/deduplikasi backup.}
\loc{Bagian~5 (Kesimpulan, arah penelitian mendatang yang diperluas).}

\reviewer{A7. Pastikan referensi [21], [41]--[44] dikutip di teks; sertakan DOI bila tersedia.}
\begin{revbox}\small
Kutip semua referensi di teks dan tambahkan DOI.
\end{revbox}
\authors{Kami meninjau seluruh referensi untuk memastikan masing-masing dikutip di teks. DOI
sedang ditambahkan untuk semua sumber yang memilikinya (mis. survei Kapoor dkk.,
\texttt{doi:10.3390/su14010008}). Referensi yang ternyata tidak memiliki DOI atau data bibliografi
lengkap yang dapat diverifikasi akan diganti dengan sumber terindeks yang memadai.}
\loc{Bagian Daftar Pustaka (entri diperbarui dan kutipan dalam teks).}

% =====================================================================
\section*{Tanggapan untuk Reviewer B}
% =====================================================================

\reviewer{B1. Abstrak harus menyatakan bahwa sistem divalidasi menggunakan model serangan
tersimulasi pada dataset sintetis berisi 43 berkas.}
\begin{revbox}\small
Nyatakan validasi dengan model tersimulasi pada dataset sintetis 43 berkas di abstrak.
\end{revbox}
\authors{Sudah dilakukan. Abstrak kini menyatakan secara eksplisit simulasi terkontrol dan dataset
sintetis 43 berkas.}
\loc{Abstrak (direvisi).}

\reviewer{B2. Pendahuluan memerlukan paragraf tentang masalah matematis efisiensi kompresi pada
data ber-entropi tinggi (terenkripsi) vs entropi rendah (\textit{plaintext}), serta \textit{research gap} yang terdefinisi.}
\begin{revbox}\small
Tambahkan paragraf tentang masalah matematis entropi/kompresi dan definisikan research gap.
\end{revbox}
\authors{Kami menambahkan paragraf berbasis entropi Shannon, $H(X)=-\sum_i p(x_i)\log_2 p(x_i)$,
yang menjelaskan batas kompresi dan mengapa data ber-entropi tinggi (terenkripsi/sudah
terkompresi) nyaris tidak terkompresi, sementara teks ber-entropi rendah terkompresi baik. Hal ini
secara langsung memotivasi koreksi urutan compress-then-encrypt dan \textit{research gap} yang terdefinisi.}
\loc{Bagian~1 (paragraf research gap/entropi).}

\reviewer{B3. Gambar 1, 2, 3, 10 harus dirujuk dan dijelaskan di teks sebelum ditampilkan; Gambar
14, 15, 16 tidak ditemukan.}
\begin{revbox}\small
Rujuk dan jelaskan semua gambar sebelum tampil; kembalikan Gambar 14--16 yang hilang.
\end{revbox}
\authors{Semua gambar kini dirujuk dan dijelaskan di teks sebelum ditampilkan. Set gambar lengkap
(termasuk gambar pemulihan/restore) telah dibangun ulang dari keluaran eksperimen yang sebenarnya
dan disertakan dalam naskah revisi.}
\loc{Bagian~3 dan~4 (seluruh rujukan gambar dan keterangannya).}

\reviewer{B4. Tabel 3, 9, 11 harus dirujuk dan dijelaskan di teks sebelum ditampilkan.}
\begin{revbox}\small
Rujuk dan jelaskan Tabel 3, 9, 11 sebelum tampil.
\end{revbox}
\authors{Semua tabel kini diperkenalkan dan dijelaskan di teks sebelum ditampilkan, termasuk tabel
dataset (dengan total eksplisit 43 berkas), tabel perubahan \textit{hash}/metadata, dan tabel
akurasi deteksi.}
\loc{Bagian~3 (tabel dataset) dan Bagian~4 (tabel hasil).}

\reviewer{B5. Kesimpulan sebaiknya lebih ringkas: rangkum temuan utama, jawab tujuan, serta nyatakan
kontribusi dan arah penelitian mendatang.}
\begin{revbox}\small
Buat kesimpulan lebih ringkas dan berorientasi tujuan.
\end{revbox}
\authors{Kesimpulan ditulis ulang untuk merangkum temuan kuantitatif utama, mengonfirmasi tujuan,
menyatakan kontribusi, dan mencantumkan arah penelitian mendatang dengan jelas.}
\loc{Bagian~5 (Kesimpulan).}

\reviewer{B6. Referensi [21] belum dikutip di teks; seluruh referensi wajib dikutip.}
\begin{revbox}\small
Pastikan referensi [21] dan semua lainnya dikutip.
\end{revbox}
\authors{Kami memastikan seluruh referensi, termasuk panduan keamanan penyimpanan NIST, kini
dikutip di teks.}
\loc{Daftar Pustaka dan kutipan dalam teks.}

% =====================================================================
\section*{Tanggapan untuk Reviewer C (nilai: 8)}
% =====================================================================

\reviewer{C1. Terdapat kesalahan tata bahasa Inggris (kehilangan artikel/subjek/kata kerja bantu),
mis. ``evaluate the resilience\ldots'' seharusnya ``To evaluate the resilience\ldots''.}
\begin{revbox}\small
Lakukan penyuntingan bahasa Inggris secara menyeluruh.
\end{revbox}
\authors{Naskah telah menjalani penyuntingan bahasa menyeluruh untuk memperbaiki tata bahasa,
menambahkan artikel/subjek yang hilang, serta meningkatkan konsistensi dan keterbacaan.}
\loc{Di sepanjang naskah.}

\reviewer{C2. Ambiguitas ukuran dataset: deskripsi menyebut 9 CSV + 9 LOG/TXT + 14 JPG + 11 XLSX =
43 berkas, tetapi teks berulang kali melaporkan restorasi 113 berkas.}
\begin{revbox}\small
Perjelas ketidaksesuaian 43 vs 113 berkas.
\end{revbox}
\authors{Hal ini telah dikoreksi sepenuhnya. Dataset adalah \textbf{43 berkas} secara konsisten.
Setiap kemunculan ``113'' sebelumnya dikoreksi menjadi 43 (abstrak, narasi, tabel RTO, dan tabel
akurasi deteksi), serta baris Total (43) ditambahkan ke tabel dataset. Eksperimen yang dijalankan
ulang mengonfirmasi pemulihan 43/43 berkas.}
\loc{Abstrak, Bagian~3 (tabel dataset), Bagian~4 (tabel RTO dan deteksi).}

\reviewer{C3. Hasil FPR 0\% dan FNR 0\% terlalu berlebihan; kesempurnaan semacam itu jarang pada
deteksi ransomware dunia nyata.}
\begin{revbox}\small
Bahas hasil FPR/FNR 0\% secara lebih hati-hati.
\end{revbox}
\authors{Kami menambahkan pembahasan eksplisit yang menegaskan bahwa FPR/FNR 0\% merupakan
\emph{konsekuensi logis} dari pemeriksaan integritas berbasis \textit{hash} yang deterministik pada
kondisi terkontrol (perubahan 1 \textit{byte} pun mengubah nilai SHA-256), dan \textbf{bukan} klaim
kesempurnaan deteksi dunia nyata. Kami menyatakan keterbatasannya terhadap penghindaran canggih
(FPE, \textit{fileless}, LotL) dan memaknai metrik ini sebagai validasi fungsional lapisan
mitigasi/pemulihan.}
\loc{Bagian~4 (paragraf diskusi FPR/FNR).}

\reviewer{C4. Deskripsikan mekanisme seleksi kompresi secara eksplisit sebagai berbasis aturan.}
\begin{revbox}\small
Nyatakan dengan jelas bahwa seleksi kompresi bersifat berbasis aturan.
\end{revbox}
\authors{Seleksi kompresi kini dideskripsikan secara eksplisit sebagai strategi \emph{berbasis
aturan} yang ditentukan oleh tipe dan ukuran berkas, dan hal ini dinyatakan di Pendahuluan, Metode,
dan Kesimpulan.}
\loc{Bagian~1, 3, dan~5.}

\reviewer{C5. Perkuat validasi eksperimen dengan analisis statistik tambahan.}
\begin{revbox}\small
Tambahkan analisis statistik.
\end{revbox}
\authors{Kami memperkuat dasar empiris dengan dua cara. Pertama, kami menjalankan ulang eksperimen
penuh pada dua platform (Windows Server 2022 dan Ubuntu 22.04 di AWS EC2) dan menambahkan tabel
konsistensi lintas-platform. Kedua, kami menambahkan analisis statistik khusus dengan \textbf{5
uji-berulang} per algoritma, melaporkan rata-rata~$\pm$~standar deviasi waktu kompresi,
\textit{throughput}, dan utilisasi CPU/RAM (tabel statistik kompresi). Pengukuran ini mengonfirmasi
stabilitas dan reproduktibilitas hasil. Demi kejelasan dan pelaporan yang jujur, kami kini
menyatakan secara eksplisit bahwa komputer Windows lokal adalah \emph{lingkungan pengembangan
awal}, sedangkan lingkungan AWS EC2 adalah \emph{lingkungan pengukuran/validasi utama} tempat
seluruh hasil kuantitatif diperoleh; konfigurasi Windows asli tetap dipertahankan (tidak dihapus)
untuk menjaga kesinambungan dengan versi yang mula-mula diajukan.}
\loc{Bagian~4 (subbagian lintas-platform, tabel konsistensi, dan subbagian analisis statistik baru).}

% =====================================================================
\section*{Persyaratan tambahan jurnal}
% =====================================================================
\authors{Kami juga telah melengkapi bagian Pendanaan, Kontribusi Penulis (CRediT), Konflik
Kepentingan, dan Ketersediaan Data, serta menyertakan biografi penulis. Revisi ini diajukan dengan
ID naskah yang sama.}

\vspace{6pt}
\noindent Kami kembali mengucapkan terima kasih kepada Editor dan para reviewer atas masukan
berharga yang telah meningkatkan kualitas naskah ini secara signifikan.

\end{document}
